%\setcounter{section}{-1}% set to Kapitel 0 % Macht Probleme, wenn zuvor schon \section{...} aufgerufen wurde z.B. bei Zusammenfügen von mehreren Modulen

\section{Wiederholung}
%\begin{frame}\ftx{\secname}
% Kapitel 1
\s{% Skript-only
    Um das frequenz- und zeitabhängige Verhalten elektrischer Netzwerke zu beschreiben, 
    werden in diesem Kapitel kurz einige wichtige Grundlagen zusammengefasst. 

    Zur Bearbeitung vorrausgesetzt werden Kenntnisse über elektrische Grundgrößen, 
    lineare passive Bauteile $R$, $L$ und $C$, Netzwerkberechnungen, 
    sowie die komplexe Wechselstromrechnung.
}
%\end{frame}

\newvideofile{Wiederholung}{Wiederholung - Frequenzabhängigkeit elektrischer Bauelemente}
\v{% Video only Frame
\begin{frame}\ftx{\secname}
    \begin{Lernziele}{Video - Wiederholung - Frequenzabhängigkeit elektrischer Bauelemente}
        Nach diesem Video werden Sie die folgenden Lernziele erreicht haben:
        \begin{itemize}\b{\setlength\itemsep{0.7em}}
            \item{die Phasenverschiebung bei $R$, $L$ und $C$ bei AC Anregung kennen,}% Wissen
            \item{das Verhalten von Induktivitäten und Kapazitäten kennen, und} % Wissen
            \item{die wichtigsten Eigenschaften von $R$, $L$ und $C$ im Vergleich verstehen.} % Können
        \end{itemize}
    \end{Lernziele}
\speech{01Wiederholung1Lernziele}{1}{%
Wilkommen zum ersten Video des Moduls Schaltungen variabler Frequenz.
In diesem Video geht es um die Wiederholung einiger Grundlagen bezüglich der Frequenzabhängigkeit elektrischer Bauelemente.
Zu den Lernzielen dieses Videos gehört, dass Sie die Phasenverschiebung bei Widerständen, Induktivitäten und Kapazitäten kennen,
das frequenzvariable Verhalten von Induktivitäten und Kapazitäten kennen,
und die wichtigsten Eigenschaften der linearen Zweipole Widerstand, Induktivität und Kapazität in Bezug auf Frequenzabhängigkeit zu verstehen.
}% Ende speech
\end{frame}
}% Ende Video only Frame

%%%%%%%%%%%%%%%%%%%%%%%%%%%%%%%%%%%%%%%%%%%%%%%%%%%%%%%%%%%%%%%%%%%%%%%%%%%%%%%%%%%%%

\subsection{Frequenzabhängigkeit elektrischer Bauelemente}
\begin{frame}\ftx{\subsecname}
\s{% Skript-only
    Rein ohmsche Widerstände $R$ können elektrische keine Energie speichern. 
    Spannungen und Ströme sind proportional zueinander und zu jedem Zeitpunkt in Phase.
    Das Verhalten ist zeit- und frequenzunabhängig.

    Induktivitäten $L$ und Kapazitäten $C$ hingegen können Energie speichern und abgeben.
    Dieser Vorgang ist inert (träge), benötigt also eine gewisse Zeit. 
    Dadurch verhalten sie sich frequenzabhängig. 

    Ströme und Spannungen stehen bei Induktivitäten und Kapazititäten in (zeitlich) differentiellem 
    linearen Verhältnis zueinander. Dadurch entsteht bei beiden Bauteiltypen eine Phasenverschiebung 
    zwischen Spannung und Strom von betragsmäßig $90\ \degree$.

    Abbildung \ref{fig:bauteile:phasenverschiebung} zeigt die Spannungs- und Stromverläufe für $R$, $L$ und $C$
    bei Erregung mit einer Wechselspannung $u_q = U \cdot \sin(\omega t)$ zum Vergleich.

    \begin{figure}[H]\centering
        \begin{subfigure}{0.68\textwidth}\centering
            \includegraphics{Tikz/pdf/plot_bauteile_phasenverschiebung_r.pdf}
        \end{subfigure}%
        \begin{subfigure}{0.28\textwidth}\centering
            \includegraphics{Tikz/pdf/circ_bauteil_R_u_i.pdf}% R
        \end{subfigure}\hfill% newline
        \begin{subfigure}{0.68\textwidth}\centering
            \includegraphics{Tikz/pdf/plot_bauteile_phasenverschiebung_l.pdf}
        \end{subfigure}%
        \begin{subfigure}{0.28\textwidth}\centering
            \includegraphics{Tikz/pdf/circ_bauteil_L_u_i.pdf}% L
        \end{subfigure}\hfill% newline
        \begin{subfigure}{0.68\textwidth}\centering
            \includegraphics{Tikz/pdf/plot_bauteile_phasenverschiebung_c.pdf}
            \caption{Strom- und Spannungsverläufe}
        \end{subfigure}%
        \begin{subfigure}{0.28\textwidth}\centering
            \includegraphics{Tikz/pdf/circ_bauteil_C_u_i.pdf}% C
            \caption{Schaltsymbole}
        \end{subfigure}
        \caption{Phasenverschiebung bei $R$, $L$ und $C$}\label{fig:bauteile:phasenverschiebung}
    \end{figure}

    Aufgrund der Phasenverschiebung bei $L$ und $C$ oszilliert deren Leistung 
    (Energieaufnahme und -abgabe), ist über eine Periode gemittelt jedoch immer null.
    Induktivitäten und Kapazitäten können daher keine Wirkleistung, sondern nur Blindleistung
    verrichten, weshalb sie auch Blindwiderstände genannt werden.
}
\b{% Folien-only
    %3 AC Plots: u(t), i(t) mit $\varphi$ eingezeichnet für $R$, $L$, $C$
    \begin{minipage}{\textwidth}\centering
        \onslide<1->{%
        \begin{minipage}{0.68\textwidth}\centering
            \resizebox{0.7\textwidth}{!}{\includegraphics{Tikz/pdf/plot_bauteile_phasenverschiebung_r.pdf}}   % plot \onslide<1->{...}        <-- OVERLAY
        \end{minipage}%
        }%
        \begin{minipage}{0.28\textwidth}\centering
            \includegraphics{Tikz/pdf/circ_bauteil_R_u_i.pdf}% R
        \end{minipage}\hfill% newline
        \onslide<2->{%
        \begin{minipage}{0.68\textwidth}\centering
            \resizebox{0.7\textwidth}{!}{\includegraphics{Tikz/pdf/plot_bauteile_phasenverschiebung_l.pdf}}   % plot \onslide<2->{...}        <-- OVERLAY
        \end{minipage}%
        }%
        \begin{minipage}{0.28\textwidth}\centering
            \includegraphics{Tikz/pdf/circ_bauteil_L_u_i.pdf}% L
        \end{minipage}\hfill% newline
        \onslide<3->{%
        \begin{minipage}{0.68\textwidth}\centering
            \resizebox{0.7\textwidth}{!}{\includegraphics{Tikz/pdf/plot_bauteile_phasenverschiebung_c.pdf}}   % plot \onslide<3->{...}        <-- OVERLAY
        \end{minipage}%
        }%
        \begin{minipage}{0.28\textwidth}\centering
            \includegraphics{Tikz/pdf/circ_bauteil_C_u_i.pdf}% C
        \end{minipage}
    \end{minipage}
    \speech{folie1}{1}{Dies ist ein Beispiel Text. A + i mal B gleich C komplex. Hundert Volt Spannung.}
}% end Folien-only
\speech{01Wiederholung2FrequenzabhaengigkeitElBauelemente}{1}{%
In diesem Abschnitt werden wir die Phasenverschiebung bei Widerständen, Induktivitäten und Kapazitäten wiederholen.
Rechts sehen wir dazu die Schaltsymbole der jeweiligen Bauelemente, links die Spannungs- und Stromverläufe bei Anregung mit einer sinusförmigen Wechselspannung.
Beginnen wir beim elektrischen Widerstand R. 
Wie im obigen Diagramm dargestellt ist, sind Spannung und Strom bei einem Widerstand immer in Phase.
Das bedeutet, dass Maxima, Minima und Nulldurchgänge von Spannung und Strom gleichzeitig auftreten.
}% Ende speech
\speech{01Wiederholung3FrequenzabhaengigkeitElBauelemente}{2}{%
Bei Induktivitäten L hingegen eilt der Strom der Spannung um 90 Grad hinterher bei Anregung mit einer sinusförmigen Wechselgröße.
Das bedeutet, dass der Strom das Maximum erst eine Viertel Periode nach der Spannung erreicht.
}% Ende speech
\speech{01Wiederholung4FrequenzabhaengigkeitElBauelemente}{3}{%
Bei Kapazitäten C ist es genau umgekehrt.
Hier eilt die Spannung dem Strom um 90 Grad hinterher.
Das bedeutet, dass die Spannung ihr Maximum erst eine Viertel Periode nach dem Strom erreicht.
Die Phasenverschiebungen bei Induktivitäten und Kapazitäten beträgt immer betragsmäßig 90 Grad.
Dadurch können diese Bauelemente keine Wirkleistung verrichten, sondern nur Blindleistung und werden daher auch Blindwiderstände genannt.
}% Ende speech
\end{frame}

%%%%%%%%%%%%%%%%%%%%%%%%%%%%%%%%%%%%%%%%%%%%%%%%%%%%%%%%%%%%%%%%%%%%%%%%%%%%%%%%%%%%%

\subsubsection{Verhalten von Induktivitäten und Kapazitäten}
\begin{frame}\ftx{\subsubsecname}
\s{% Skript-only
    Induktivitäten als ideale Bauteile speichern durch den Effekt der (Selbst-)Induktion Energie im Magnetfeld.
    Kapazitäten als ideale Bauteile hingegen speichern Energie im elektrischen Feld. 
    Beschrieben werden beide Effekte durch das Induktion- beziehungsweise das Gaußsche Gesetz.

    Tabelle \ref{tab:vergleich:induktivitaetkapazitaet} listet die wichtigsten Unterschiede 
    im Verhalten von Induktivitäten und Kapazitäten qulitativ als Übersicht auf.

    \begin{table*}[h]\centering
    \caption{Vergleich von Induktivität und Kapazität im Verhalten}
    \label{tab:vergleich:induktivitaetkapazitaet}
    \begin{tabular}{ccc}
        \toprule&\textbf{Induktivität} &\textbf{Kapazität}\\
        \midrule
        \textbf{Gesetz}         &Induktionsgesetz       &Gaußsches Gesetz   \vphantom{$\Big|$}\\
        \textbf{Energiespeich.} &im Magnetfeld          &im Elektr. Feld    \vphantom{$\Big|$}\\
        \textbf{stetig}         &Strom                  &Spannung           \vphantom{$\Big|$}\\
        \textbf{bei Gleichspg.} &Kurzschluss            &offen              \vphantom{$\Big|$}\\
        \textbf{bei Hochfreq.}  &offen                  &Kurzschluss        \vphantom{$\Big|$}\\
        \bottomrule
    \end{tabular}
    \end{table*}

    Die (Selbst-)Induktivität $L$ als Eigenschaft kann vereinfacht als \glqq Trägheit \grqq des Stroms
    betrachtet werden. Ströme in Induktivitäten sind stetig und eilen der Spannung hinterher.
    Bei hohen Frequenzen sperrt die Induktivität, bei Gleichspannung verhält sie sich wie ein Kurzschluss.
    
    Die Kapazität $C$ kann dem gegenüber vereinfacht als \glqq Trägheit \grqq der Spannung betrachtet werden.
    Spannungen in Kapazitäten sind stetig und eilen dem Strom hinterher. 
    Bei Gleichstrom sperrt die Kapazität, bei hohen Frequenzen verhält sie sich wie ein Kurzschluss.

    \begin{figure}[H]\centering
        \begin{subfigure}{0.45\textwidth}\centering
            \includegraphics[width=0.75\textwidth]{./Bilder/Spule.png}
        \end{subfigure}%
        \begin{subfigure}{0.45\textwidth}\centering
            \includegraphics[width=0.75\textwidth]{./Bilder/Kondensator.png}
        \end{subfigure}
        \caption{Reale Induktivät (Spule, links) und Kapazitäten (Kondensator, rechts)}\label{fig:bauteile:realebauteile}
    \end{figure}
    Abbildung \ref{fig:bauteile:realebauteile} zeigt reale Bauteile, die Induktivitäten und Kapazitäten realisieren. 
    Gezeigt sind eine Spule (Induktivität) und ein Kondensator (Kapazität).
}
\b{% Folien-only
\begin{columns}
\column{0.4\textwidth}
    \textbf{Induktivität}%
    \begin{itemize}
        \item Impedanz: $\underline{Z}_{\mathrm L} = \mathrm{j}\omega L$
        \item Zeitb.: $\ecv{u_{\mathrm L}} = L \cdot \frac{\d}{\d t}\, \eci{i_{\mathrm L}}$
        \item Energie im Magnetfeld
        \item Ströme stetig (keine Sprünge)
    \end{itemize}

    \textbf{Kapazität}%
    \begin{itemize}
        \item Impedanz: $\underline{Z}_{\mathrm C} = -\mathrm{j}\frac{1}{\omega C}$
        \item Zeitb.: $\eci{i_{\mathrm C}} = C \cdot \frac{\d}{\d t}\, \ecv{u_{\mathrm C}}$
        \item Energie im elektrischen Feld
        \item Spannungen stetig (keine Sprünge)
    \end{itemize}

\column{0.15\textwidth}\centering
    \includegraphics{Tikz/pdf/circ_bauteil_L_u_i_mini.pdf}% Schaltsymbol
    \vspace{2cm}
    \includegraphics{Tikz/pdf/circ_bauteil_C_u_i_mini.pdf}% Schaltsymbol

\column{0.45\textwidth}
    \begin{figure}[H]\centering
        \includegraphics[width=0.5\textwidth]{./Bilder/Spule.png}
        \par Spule
    \end{figure}
    \begin{figure}\centering
        \includegraphics[width=0.5\textwidth]{./Bilder/Kondensator.png}
        %\includegraphics[width=0.5\textwidth]{./Bilder/Kondensator_cut.png}
        \par Kondensator
    \end{figure}
\end{columns}
}% end Folien-only
\speech{01Wiederholung3InduktivitaetenKapazitaeten}{1}{%
Betrachten wir nun das Verhalten von Induktivitäten und Kapazitäten im Vergleich.
Von links nach rechts sehen Sie zunächst die wichtigsten Eigenschaften von Induktivitäten und Kapazitäten,
in der Mitte die jeweiligen Schaltsymbole und rechts Fotos von realen Bauelementen, mit denen Induktivitäten und Kapazitäten realisiert werden.
Beginnen wir mit der Induktivität.
}% Speech Ende
\silence{1}%
\speech{01Wiederholung3InduktivitaetenKapazitaeten}{1}{%
In der Wechselstromrechnung wird die Induktivität durch ihre Impedanz beschrieben, die proportional zur Kreisfrequenz und Induktivität ist.
Im Zeitbereich ausgedrückt, ist die U proportional zu L mal der zeitlichen Ableitung von I.
Dadurch dass die Energie bei Induktivitäten im Magnetfeld gespeichert wird, sind Ströme in Induktivitäten stetig, das heißt sie können keine Sprünge aufweisen.
Bei Kapazitäten verhält es sich genau umgekehrt. 
Die Admittanz ist proportional zur Kreisfrequenz und Kapazität, dadurch ist die Impedanz negativ und umgekehrt proportional zur Kreisfrequenz und Kapazität.
Im Zeitbereich ausgedrückt, ist der I proportional zu C mal der zeitlichen Ableitung von U.
Dadurch dass die Energie bei Kapazitäten im elektrischen Feld gespeichert wird, sind Spannungen in Kapazitäten stetig, das heißt sie können keine Sprünge aufweisen.
}% Ende speech
\end{frame}

%%%%%%%%%%%%%%%%%%%%%%%%%%%%%%%%%%%%%%%%%%%%%%%%%%%%%%%%%%%%%%%%%%%%%%%%%%%%%%%%%%%%%

\subsubsection{Vergleich der linearen Zweipole $R$, $L$ und $C$}
\begin{frame}\ftx{\subsubsecname}

    \s{\ref{tab:vgl:bauteile:rlc} zeigt eine Gegenüberstellung der Größen $R$, $L$ und $C$ und fasst deren wichtigsten Eigenschaften zusammen. }

    \begin{table}[h]\centering%
    \s{\caption{Vergleich lineare Bauteile $R$, $L$, $C$}\label{tab:vgl:bauteile:rlc}}%
    \resizebox{\textwidth}{!}{% for beamer page fit
    \begin{tabular}{ |c|c|c|c|c| }
        \hline &&&&\\[-6pt]
        \textbf{Größe}                          % Col 1
                & \textbf{Allgemein}            % Col 2
                & \textbf{El. Widerstand}       % Col 3
                & \textbf{Induktivität}         % Col 4
                & \textbf{Kapazität} \\[+4pt]   % Col 5
        \hline%
        %\textbf{Symbol}
        \includegraphics{Tikz/pdf/circ_bauteil_tabelle_text_symbol.pdf}% Symbol (vertikal zentriert); Schriftart entspricht Textfont im Skript (mit Serifen), nicht in den Folien (ohne Serifen)
                & \includegraphics{Tikz/pdf/circ_bauteil_tabelle_zweipol.pdf}
                & \includegraphics{Tikz/pdf/circ_bauteil_tabelle_R.pdf}
                & \includegraphics{Tikz/pdf/circ_bauteil_tabelle_L.pdf}
                & \includegraphics{Tikz/pdf/circ_bauteil_tabelle_C.pdf} \\[+4pt]
        \hline &&&&\\[-6pt] % Dummyline für vertikalen Abstand
        \textbf{Einheit}
                & $ \left[\textit{Form.z.}\right] = \mathrm{Einheit} $ 
                & $ \left[R\right] = \Omega\ \text{(Ohm)} $     
                & $ \left[L\right] = \mathrm{H}\ \text{(Henry)} $ 
                & $ \left[C\right] = \mathrm{F}\ \text{(Farad)} $ \\[+10pt]
        \hline %&&&&\\[-6pt]
        \textbf{Zeitbereich} $\vphantom{\bigg|}$ % vphantom für vertikalen Abstand
                & $ \frac{\d}{\d t} $ bzw. $ \int \d t $ 
                & $ \ecv{u_{\mathrm{R}}} = R \cdot \eci{i_{\mathrm R}} $     
                & $ \ecv{u_{\mathrm{L}}} = L \cdot \frac{\d}{\d t}\, \eci{i_{\mathrm{L}}} $ 
                & $ \eci{i_{\mathrm{C}}} = C \cdot \frac{\d}{\d t}\, \ecv{u_{\mathrm{C}}} $ \\[+10pt]
        \hline %&&&&\\[-6pt]
        \textbf{Frequenzb.} $\vphantom{\bigg|}$ % vphantom für vertikalen Abstand
                & $ \mathrm{j}\omega $ bzw. $ \frac{1}{\mathrm{j}\omega} $
                & $ \ecv{\underline{U}_{\mathrm R}} = R                   \cdot \eci{\underline{I}_{\mathrm R}} $ 
                & $ \ecv{\underline{U}_{\mathrm L}} = \mathrm{j}\omega L  \cdot \eci{\underline{I}_{\mathrm L}} $
                & $ \eci{\underline{I}_{\mathrm C}} = \mathrm{j}\omega C  \cdot \ecv{\underline{U}_{\mathrm C}} $ \\[+10pt]
        \hline %&&&&\\[-6pt]
        \textbf{Impedanz} $\vphantom{\bigg|}$ % vphantom für vertikalen Abstand
                & $ \underline{Z} = \frac{\ecv{\underline{U}}}{\eci{\underline{I}}} $
                & $ \underline{Z}_{\mathrm R} = R $ 
                & $ \underline{Z}_{\mathrm L} = \mathrm{j}\omega L $
                & $ \underline{Z}_{\mathrm C} = -\mathrm{j}\frac{1}{\omega C} $ \\[+10pt]
        \hline %&&&&\\[-6pt]
        \textbf{Wirkanteil} $\vphantom{\bigg|}$ % vphantom für vertikalen Abstand
                & $ R = \Re\{\underline{Z}\} $
                & $ R = \frac{U}{I} $
                & $ 0 $
                & $ 0 $ \\[+10pt]
                \hline %&&&&\\[-6pt]
        \textbf{Blindanteil} $\vphantom{\bigg|}$ % vphantom für vertikalen Abstand
                & $ X = \Im\{\underline{Z}\} $
                & $ 0 $
                & $ X_{\mathrm L} = \omega L $
                & $ X_{\mathrm C} = -\frac{1}{\omega C}$ \\[+10pt]
        \hline
    \end{tabular}
    } % end resizebox
    \end{table}
    
% Für Schaltungen variabler Frequenz: Frequenzvariablen Eigenschaften von Bauteilen, dafür zeitvariante 
\speech{01Wiederholung4VergleichLinearerZweipoleRLC}{1}{%
Hier sehen Sie zur Zusammenfassung und Übersicht eine Gegenüberstellung der linearen Bauelemente Widerstand R, Induktivität L und Kapazität C.
Von oben nach unten sind die wichtigsten Größen und Eigenschaften der Bauelemente aufgelistet, von links nach rechts die jeweiligen Bauelemente.
Für uns relevant für die folgenden Kapitel sind vor allem die komplexen Strom- und Spannungsbeziehungen im Frequenzbereich, also die Impedanzen respektive Admittanzen.
}% Speech Ende
\silence{1}%
\speech{01Wiederholung4VergleichLinearerZweipoleRLC}{1}{%
Achten Sie bei der Angabe der Reaktanzen X, also der Imaginärteile der Impedanzen Z immer auf die Vorzeichen.
Die Reaktanz der Kapazität ist negativ, die der Induktivität positiv.
Umgekehrt ist die Suszeptanz B, das ist der Imaginärteil der Admittanz Y, bei Kapazitäten positiv und bei Induktivitäten negativ.
Bei Schwierigkeiten bei der komplexen Wechselstromrechnung empfehle ich Ihnen, das Modul 7 periodische Größen nochmals durchzuarbeiten.
Die Grundlagen der komplexen Wechselstromrechnung sind Vorraussetzung für das Verständnis der folgenden Kapitel.
Vielen Dank fürs Zuhören.
}% Speech Ende
\end{frame}

%%%%%%%%%%%%%%%%%%%%%%%%%%%%%%%%%%%%%%%%%%%%%%%%%%%%%%%%%%%%%%%%%%%%%%%%%%%%%%%%%%%

% Beschreibe in Textform wie das Verhalten von Eingangs- zu Ausgangssignalen eines Vierpols 
% mithilfe des Frequenzgangs (bzw. des Amplituden- und Phasengangs) beschrieben werden kann
\s{
\subsection{Zweitore (Vierpole)}
\begin{frame}\ftx{\subsecname}

    % Wiederholung Eintor
    Die Frequenzabhängigkeit von Eintoren (Zweipole) haben wir in Modul 3 bereits untersucht.  
    Mithilfe der komplexen Wechselstromrechnung konnten wir den frequenzabhängigen Wechselstromwiderstand 
    (Impedanz) $\underline{Z}$ einzelner Eintore bestimmen. Durch Anwendung der komplexen Strom- 
    und Spannungsteilerregel konnten wir auch die Gesamt-Impedanz linearer zweipoliger Netzwerke bestimmen. 
    Das heißt für solche, welche nur aus linearen Bauteilen wie $R$, $L$ und $C$ bestehen. 

    % Tikzpicture: Vergleich Vierpol und Zweipol
    %
    % Zeichne links einen senkrechten Zweipol mit Spannungs- (U) und Strompfeil (I) jeweils von oben nach unten
    % Zeichne rechts einen Vierpol (Eingangspole links, Ausgangspole rechts) (U1 links oben nach unten, U2 rechts oben nach unten)
    \begin{figure}[h]
        \begin{subfigure}{0.48\textwidth}\centering
            \includegraphics{Tikz/pdf/circ_eintor.pdf}
        \end{subfigure}\hfill
        \begin{subfigure}{0.48\textwidth}\centering
            \includegraphics{Tikz/pdf/circ_zweitor.pdf}
        \end{subfigure}
        \caption{Vergleich: Schaltzeichen eines allgemeinen Zwei- und Vierpols}
        \label{fig:vgl:zweipolvierpol}
    \end{figure}

    % Einführung Zweitor, wozu
    Zweitore (Vierpole) sind eine Erweiterung von Eintoren (Zweipolen).
    Sie verfügen über eine Eingangs- und eine Ausgangsseite mit jeweiligen Eingangs- und Ausgangsgrößen.  
    Das Zweitormodell eignet sich gut zur Beschreibung von (frequenzabhängigen) Übertragungseigenschaften 
    von elektrischen Netzwerken. 

    \fu{\includegraphics{Tikz/pdf/circ_zweitor_quelle_last.pdf}
    }{Verbindung einer Last über ein Zweitor mit der Quelle.\label{circ:quellezweitorlast}}% im Foliensatz im nächsten Kapitel (Freuqenzgang, Amplitudengang, Phasengang)
    
    Im einfachsten Fall wird wie in \ref{circ:quellezweitorlast} gezeigt, eine Quelle (Eintor) über ein Zweitor 
    mit einer Last (Eintor) verbunden. 
    Typische Zweitore sind beispielsweise Verstärker, Filter oder auch Transformatoren. 
\end{frame}
}% end skript only


